\documentclass[10pt,a4paper]{amsart}
\usepackage[margin=19mm]{geometry}
\usepackage{amsmath,amssymb,mathtools}
\usepackage[scr=rsfs,cal=euler]{mathalfa}
\usepackage{xfrac}
\usepackage[unicode,hidelinks,bookmarks=false]{hyperref}
\newcommand{\ie}{i.e.\ }
\newcommand{\cf}{cf.\ }
\newcommand{\resp}{respectively\ }
\newcommand{\Subsection}{\subsection}
\providecommand{\nobreakdash}{\nobreak\hbox{-}}
\setlength{\emergencystretch}{2em}
\multlinegap0pt
\usepackage{amssymb}
\newtheorem{lemma}{Lemma}
\newcommand{\Si}{\Sigma}
\renewcommand{\a}{\alpha}
\newcommand{\be}{\begin{equation}}
\newcommand{\cB}{{\mathcal B}}
\newcommand{\cC}{{\mathcal C}}
\newcommand{\cI}{\mathcal I} 
\newcommand{\ee}{\end{equation}}
\newcommand{\f}{\varphi}
\title{Well-Posed Geometric Boundary data in General Relativity, II: Twisted Dirichlet boundary data}
\author{Zhongshan An, Michael T. Anderson}
\date{Source: arXiv:2507.15567v3 (CC BY 4.0)}
\begin{document}
\maketitle

\noindent\emph{Source and licence.} Well-Posed Geometric Boundary data in General Relativity, II: Twisted Dirichlet boundary data. Version arXiv:2507.15567v3 (CC BY 4.0), distributed by arXiv under the Creative Commons Attribution 4.0 International licence, http://creativecommons.org/licenses/by/4.0/, as recorded in the arXiv licence metadata for this exact version and shown on the abstract page https://arxiv.org/abs/2507.15567v3. Full licence notice: \texttt{licenses/arxiv-2507.15567-CC-BY-4.0.txt}.\par\smallskip

\noindent\textit{Editorial scope.} Counted unit: the article's lemma corner (source0.tex lines 670--681). The article's complete original proof of it is delivered with it (source0.tex lines 683--687, one \texttt{proof} environment of 5 lines). No mathematics is written, repaired or re-derived: the statement and the proof are the article's own lines, in the article's own notation, and no retained line is substituted or re-flowed.  No further source-local context is needed: every cross-reference of every retained line resolves inside the two delivered spans.  Nothing was omitted from any retained span, so the package carries no omission record.  No retained line contains a citation, so the wrapper carries no bibliography and no bibliography entry.  No retained line contains a figure environment, an image-inclusion command or drawing code, so no image is delivered and the package declares no asset dependency.  Editorial formatting only, in addition: an AMS \texttt{amsart} preamble that restates the article's own declarations and macros used by the retained lines, namely the environments \texttt{lemma} on separate counters, together with the standard packages those lines need.  The retained environments print in this document as Lemma 1 in order of appearance, whereas the article prints the same items with the numbering of its own full document.  The complete native source of the article as distributed in the arXiv e-print for this exact version is bundled unchanged as \texttt{sources/181952/source0.tex}.
\begin{lemma} \label{corner}
In local coordinates $\{x^{\mu}\}$ near $\Si$ with $g_{11} = f^2$ and $g_{1a} = 0$, $a = 2, \dots, n$ (zero shift), one has 
\be \label{dvg}
dv_g = \sqrt{f^2 det(-g_{\cC}) + \a^2 det(g_{\Si})} dx^0 \wedge \cdots \wedge x^n,
\ee
where $\a = -g_{01}$. Consequently, 
\be \label{mu1}
\mu_g = \frac{\sqrt{f^2 det(-g_{\cC}) + \a^2 det(g_{\Si})}}{[det(-g_{\cC})]^{1/n}},
\ee
and so at the corner $\Si$, $\a$ is uniquely determined, up to sign, by the initial and boundary data in $\cI \times \cB$. 

\end{lemma}

\begin{proof}
The formula \eqref{dvg} is a simple calculation in local coordinates, as is the relation $\a = -g_{01}$ and the formula \eqref{mu1}. As noted above, 
at $\Si$, both $f^2 = g_{11}$ and the conformal factor $\f$ are determined by initial data. Hence, $g_{\cC}$ and $g_{\Si}$ are determined at 
$\Si$ by the target data in $\cI \times \cB$. Thus the second statement follows from the formula \eqref{mu1}. 
\end{proof}

\end{document}
